% Ad Atticum 16.11 --- headnote
% ad-atticum-16-11, 5 November 44 BC, in Puteolano

\textit{Cicero to Atticus, written at the Puteolan villa on
the Nones of November (5 November) 44 BC --- Perseus
dateline \textit{Scr. in Puteolano Non. Nov. a. 710 (44)}.
The longest letter of this November cluster, and politically
the most consequential. Cicero is still on the coast,
catching the eastward post each morning, and the letter
answers two from Atticus at once --- one from the Kalends,
the other from the day before.}

\textit{The opening is literary. ``Nostrum opus'' is the
Second Philippic, which Cicero has finished in draft and
sent to Atticus for marking up; the ``little vermilion-wax
marks'' are Atticus's annotations, and Cicero has been
dreading them. Atticus has approved, even culled
\textit{anthē} --- ``flowers,'' purple passages --- from it;
Cicero will soften the personal attack on the cuckolded
Antony (the gibe about Fadia, made without ``the Lucilian
phallus'') so as not to wound Sicca and Septimia, but he
will not retract the politics. He longs for ``the day when
that speech shall roam at large so freely that it makes its
way even into Sicca's house''; for the moment it is a
pamphlet circulating only among trusted hands, to be read to
Sextus Peducaeus, with Calenus and Calvena (Matius) kept
out.}

\textit{The middle of the letter turns to the philosophical
work that will become \textit{De Officiis}. Cicero is using
Panaetius as his base, has finished two books on the topic
Panaetius treated, and is waiting on Athenodorus Calvus to
send him the headings of Posidonius's treatment of the third
question --- the conflict of the honourable with the useful
--- which Panaetius announced but never wrote. The title is
fixed: \textit{kathēkon} renders ``officium,'' and the
fuller title is \textit{De Officiis}; the dedication will be
to his son.}

\textit{Section 6 is the heart of the letter. Octavian is
writing to Cicero daily, calling him to Capua, urging him to
``save the republic a second time,'' demanding he come to
Rome at once. Cicero quotes Iliad 7.93 --- ``they were
ashamed to refuse'' --- on his own ambivalence. He sees what
Octavian is doing clearly: the boy is acting with real
energy, the country towns adore him (a remarkable welcome at
Cales, an exhortation at Teanum), the march into Samnium is
under way --- but ``he is plainly a boy,'' he supposes a
senate will simply assemble when he asks, and Cicero asks
the obvious question: who will come? On account of all this,
Cicero has resolved to be at Rome sooner than planned. The
closing sections shift back to small business: the bond on
the day before the Ides, letters to be carried to the
Sicilian cities, the Lepidian holidays, Quintus's affection
for someone Atticus likes less, and a kiss for the bright
small Attica.}
